\documentclass[11pt,a4paper]{article}
\usepackage[T1]{fontenc}
\usepackage[utf8]{inputenc}
\usepackage[french]{babel}
\usepackage{mathptmx}
\usepackage[a4paper,margin=2.4cm]{geometry}
\usepackage{amsmath,amssymb}
\usepackage{booktabs,longtable,array}
\usepackage{enumitem}
\usepackage[hidelinks]{hyperref}
\setlength{\parskip}{4pt}
\setlength{\parindent}{0pt}
\setlength{\emergencystretch}{2em}
\newcommand{\bv}[1]{\mathbf{#1}}
\newcommand{\Jm}{\boldsymbol{\mathcal{J}}}
\newenvironment{expl}{\begin{itemize}[leftmargin=1.2em,itemsep=2pt,topsep=2pt]}{\end{itemize}}
\newcommand{\sens}{\item \textbf{Sens.}~}
\newcommand{\termes}{\item \textbf{Les termes.}~}
\newcommand{\origine}{\item \textbf{D'où elle vient.}~}
\newcommand{\num}{\item \textbf{Valeurs numériques.}~}
\newcommand{\role}{\item \textbf{À quoi elle sert.}~}

\begin{document}

\begin{center}
{\Large\bfseries Guide de lecture de l'article\par}
\vspace{4pt}
{\large\itshape Hybrid Flatness- and Passivity-Based Control with Adaptive RBF Compensation\\
for a 2~MW Direct-Drive PMSG Wind Energy Conversion System\par}
\vspace{6pt}
{\normalsize Abréviations, symboles et explication détaillée de chaque équation\par}
\end{center}

\tableofcontents

%======================================================================
\section{L'idée de l'article en une page}
Le système étudié est une éolienne de 2~MW à entraînement direct : la turbine fait tourner directement
une machine synchrone à aimants permanents (PMSG), sans multiplicateur. La puissance passe par deux
convertisseurs dos à dos : le convertisseur côté machine (MSC) redresse la tension du générateur, le bus
continu (\emph{DC link}) découple les deux convertisseurs, et le convertisseur côté réseau (GSC) injecte la
puissance dans le réseau 690~V, 50~Hz à travers un filtre $L$.

La commande proposée a trois couches, chacune avec un rôle précis :
\begin{enumerate}[itemsep=2pt]
\item \textbf{La platitude (FBC)} génère les références. À partir de la vitesse désirée de la turbine, elle
calcule directement le couple et le courant $i_{sq}^*$ nécessaires ; à partir de l'énergie désirée du bus
continu, elle calcule directement la puissance et le courant $i_{gd}^*$ à injecter dans le réseau.
\item \textbf{La passivité (PBC)} fait suivre ces références de courant par les deux convertisseurs. Elle
garde la structure énergétique de la machine et du filtre et ajoute de l'amortissement.
\item \textbf{Le réseau de neurones RBF adaptatif} apprend en ligne l'erreur du modèle du couple
aérodynamique (et des paramètres) et la compense dans la loi de vitesse, sans changer les gains.
\end{enumerate}
L'article montre ensuite, par une analyse de Lyapunov, que tous les signaux restent bornés, puis vérifie
par simulation (avec bruit de mesure et harmoniques du réseau) que les grandeurs suivent leurs références,
et que le réseau de neurones réduit l'erreur de vitesse de 26.5 à 44.8\,\%.

\textbf{Organisation de l'article.} Section~2 : modèle du système. Section~3 : objectifs de commande et
MPPT. Section~4 : platitude. Section~5 : passivité. Section~6 : réseau RBF. Section~7 : stabilité.
Section~8 : simulations. Section~9 : conclusion.

%======================================================================
\section{Abréviations}
\begin{longtable}{@{}p{1.9cm}p{5.6cm}p{7.4cm}@{}}
\toprule
Abréviation & En anglais & Signification\\
\midrule
\endhead
WECS & Wind Energy Conversion System & Système de conversion de l'énergie éolienne (turbine + générateur + convertisseurs + réseau).\\
PMSG & Permanent Magnet Synchronous Generator & Générateur synchrone à aimants permanents.\\
D-PMSG & Direct-drive PMSG & PMSG à entraînement direct (sans multiplicateur).\\
MSC & Machine-Side Converter & Convertisseur côté machine (redresseur MLI qui commande les courants du générateur).\\
GSC & Grid-Side Converter & Convertisseur côté réseau (onduleur MLI qui injecte la puissance au réseau).\\
DC & Direct Current & Continu ; \emph{DC link} = bus continu entre les deux convertisseurs.\\
AC & Alternating Current & Alternatif.\\
PWM & Pulse Width Modulation & Modulation de largeur d'impulsion (MLI).\\
SVPWM & Space Vector PWM & MLI vectorielle ; limite linéaire de la tension : $V_{dc}/\sqrt3$.\\
FOC & Field-Oriented Control & Commande vectorielle classique (boucles PI dans le repère $dq$).\\
PI & Proportional-Integral & Correcteur proportionnel-intégral.\\
FBC & Flatness-Based Control & Commande par platitude différentielle.\\
PBC & Passivity-Based Control & Commande basée sur la passivité (mise en forme de l'énergie + injection d'amortissement).\\
IDA-PBC & Interconnection and Damping Assignment PBC & Variante de la PBC qui impose l'interconnexion et l'amortissement (réf. [20]).\\
PCH & Port-Controlled Hamiltonian & Système hamiltonien commandé par ports (forme énergétique d'un modèle).\\
SMC & Sliding Mode Control & Commande par modes glissants.\\
MPC & Model Predictive Control & Commande prédictive.\\
RBF & Radial Basis Function & Fonction à base radiale (ici gaussienne).\\
RBFNN & RBF Neural Network & Réseau de neurones à fonctions de base radiales.\\
ANN & Artificial Neural Network & Réseau de neurones artificiels (ici, le réseau RBF).\\
MPPT & Maximum Power Point Tracking & Recherche du point de puissance maximale.\\
TSR & Tip-Speed Ratio & Vitesse spécifique $\lambda=\Omega R/v$ ; la MPPT \og TSR\fg{} impose $\lambda=\lambda^\star$.\\
OTC & Optimal Torque Control & MPPT par couple optimal $T_e^*=K_{opt}\Omega^2$ (non utilisée dans l'article).\\
PLL & Phase-Locked Loop & Boucle à verrouillage de phase (synchronisation au réseau ; supposée idéale ici).\\
THD & Total Harmonic Distortion & Taux de distorsion harmonique du courant.\\
TDD & Total Demand Distortion & Distorsion rapportée au courant nominal (norme IEEE 519).\\
IEEE & Institute of Electrical and Electronics Engineers & Ici : norme IEEE Std 519-2014 sur les harmoniques.\\
PF & Power Factor & Facteur de puissance ($Q_g=0$ $\Rightarrow$ PF $=1$).\\
RK4 & Runge--Kutta (order 4) & Méthode d'intégration numérique d'ordre 4 à pas fixe.\\
ZOH & Zero-Order Hold & Bloqueur d'ordre zéro (tension maintenue constante pendant $T_s$).\\
UUB & Uniformly Ultimately Bounded & Uniformément finalement borné : l'erreur entre, après un temps fini, dans un ensemble borné et y reste.\\
IAE & Integral of Absolute Error & Intégrale de la valeur absolue de l'erreur.\\
ISE & Integral of Squared Error & Intégrale du carré de l'erreur.\\
RMSE & Root-Mean-Square Error & Erreur quadratique moyenne (racine).\\
MAE & Mean Absolute Error & Erreur absolue moyenne.\\
std & standard deviation & Écart type.\\
rms & root mean square & Valeur efficace.\\
LCL & Inductor--Capacitor--Inductor & Filtre $LCL$ (perspective de l'article).\\
HIL & Hardware-In-the-Loop & Validation \og matériel dans la boucle\fg{} (perspective).\\
$dq$ / $abc$ & -- & Repère tournant de Park / grandeurs triphasées.\\
\bottomrule
\end{longtable}

%======================================================================
\section{Symboles}
Les valeurs sont celles de l'article (Tables 1 et 3). \og Calc.\fg{} = valeur calculée.
\begin{longtable}{@{}p{2.2cm}p{6.9cm}p{1.9cm}p{3.9cm}@{}}
\toprule
Symbole & Signification & Unité & Valeur\\
\midrule
\endhead
\multicolumn{4}{@{}l}{\emph{Turbine et vent}}\\
$v$, $v_m$ & vitesse du vent ; vent mesuré par l'anémomètre & m/s & 6--11\\
$\rho$ & masse volumique de l'air & kg/m$^3$ & 1.225\\
$R$ & rayon du rotor & m & 40\\
$\Omega$ & vitesse mécanique du rotor & rad/s & $\Omega_n=2.23$\\
$\lambda$, $\lambda^\star$ & vitesse spécifique, valeur optimale & -- & $\lambda^\star=8.1$\\
$\beta$ & angle de calage des pales (pitch) & deg & 0\\
$C_p$, $C_{p,\max}$ & coefficient de puissance, maximum & -- & 0.48\\
$P_w$, $T_w$ & puissance et couple aérodynamiques & W, N$\cdot$m & --\\
$\hat T_w$ & couple aérodynamique estimé par le modèle & N$\cdot$m & --\\
$P_n$, $v_n$, $T_n$ & puissance, vent et couple nominaux & W, m/s, N$\cdot$m & 2~MW, 11.0, 0.897~MN$\cdot$m\\
$J$, $B$ & inertie totale, frottement visqueux & kg$\cdot$m$^2$, N$\cdot$m$\cdot$s & $2.8\times10^6$, $10^3$\\
\multicolumn{4}{@{}l}{\emph{Générateur}}\\
$p$ & nombre de paires de pôles & -- & 26\\
$R_s$, $L_s$ & résistance et inductance statoriques ($L_d=L_q=L_s$) & $\Omega$, H & 0.78~m$\Omega$, 1.57~mH\\
$\psi_f$ & flux des aimants & Wb & 9.18\\
$\omega_e$ & pulsation électrique $=p\Omega$ & rad/s & 58 à $\Omega_n$\\
$v_{sd}, v_{sq}$; $\bv v_s$ & tensions statoriques ; vecteur $[v_{sd}\;v_{sq}]^T$ & V & --\\
$i_{sd}, i_{sq}$; $\bv i_s$ & courants statoriques ; vecteur & A & $i_{sq}\approx-2.5$~kA à $T_n$\\
$T_e$ & couple électromagnétique ($<0$ en générateur) & N$\cdot$m & --\\
$\kappa$ & constante de couple $\frac32p\psi_f$ & N$\cdot$m/A & 358 (calc.)\\
\multicolumn{4}{@{}l}{\emph{Bus continu, filtre, réseau}}\\
$C$, $V_{dc}$, $V_{dc}^*$ & capacité, tension du bus et sa référence & F, V & 20~mF, 1200~V\\
$P_{dc}$, $P_i$ & puissance entrant dans le bus (du MSC), sortant vers le GSC & W & --\\
$\bv v_i$ & tension du GSC & V & --\\
$L_f$, $R_f$ & inductance et résistance du filtre & H, $\Omega$ & 0.15~mH, 2~m$\Omega$\\
$\bv i_g=[i_{gd}\;i_{gq}]^T$ & courant injecté au réseau & A & 2.37~kA crête nominal\\
$\bv v_g=[v_{gd}\;v_{gq}]^T$ & tension du réseau ($v_{gq}=0$) & V & $v_{gd}=563$~V crête\\
$\omega_g$ & pulsation du réseau & rad/s & $2\pi\cdot50$\\
$P_g$, $Q_g$ & puissances active et réactive injectées & W, var & --\\
$E$ & énergie totale stockée dans le système & J & --\\
$y_2$ & énergie du bus et du filtre (sortie plate) & J & 14.4~kJ (bus)\\
\multicolumn{4}{@{}l}{\emph{Commande}}\\
$\Omega_{opt}$, $\Omega^*$ & vitesse optimale brute ; référence filtrée & rad/s & --\\
$\omega_r$ & pulsation du filtre de référence & rad/s & 1\\
$e_\Omega$, $e_y$ & erreurs de vitesse et d'énergie (mesure moins référence) & rad/s, J & --\\
$\nu_1$ & accélération imposée (nouvelle entrée) & rad/s$^2$ & --\\
$k_{11},k_{12}$ & gains de la boucle de vitesse & s$^{-1}$, s$^{-2}$ & 4, 4 (0 avec le RBF)\\
$k_{21},k_{22}$ & gains de la boucle du bus continu & s$^{-1}$, s$^{-2}$ & 120, 3600\\
$\bv e_s$, $\bv e_g$ & erreurs de courant MSC et GSC & A & --\\
$\bv z_s$, $\bv z_g$ & intégrales des erreurs de courant & A$\cdot$s & --\\
$R_a$, $K_i$ & amortissement et gain intégral PBC (MSC) & $\Omega$, $\Omega$/s & 1.57, 393\\
$R_b$, $K_{ig}$ & amortissement et gain intégral PBC (GSC) & $\Omega$, $\Omega$/s & 0.148, 37.5\\
$\omega_\Omega,\omega_v,\omega_i$ & pulsations des pôles doubles des boucles & rad/s & 2, 60, 500\\
$\Jm$ & matrice antisymétrique $\bigl[\begin{smallmatrix}0&1\\-1&0\end{smallmatrix}\bigr]$ & -- & --\\
$\bv e_q$ & vecteur $[0\;1]^T$ (axe $q$) & -- & --\\
$H_s,H_e,H_g$ & énergies (fonctions de stockage) & J & --\\
$T_s$ & période d'échantillonnage $1/(2f_{sw})$ & s & 231.5~$\mu$s\\
$f_{sw}$ & fréquence de commutation & Hz & 2160\\
\multicolumn{4}{@{}l}{\emph{Réseau de neurones et stabilité}}\\
$\Delta$, $\hat\Delta$ & incertitude globale $T_w-\hat T_w$ et son estimation & N$\cdot$m & $\approx 83$~kN$\cdot$m (cas 3)\\
$\bv z$ & entrée du réseau $[\Omega/\Omega_n,\;v_m/v_n]^T$ & -- & $\approx$0.5--1.1\\
$\Phi_j$, $\bv c_j$, $b$ & gaussiennes, centres, largeur & -- & 25 unités, $b=0.15$\\
$\bv W^*$, $\hat{\bv W}$, $\tilde{\bv W}$ & poids idéaux, estimés, erreur $\bv W^*-\hat{\bv W}$ & N$\cdot$m & --\\
$\varepsilon$, $\bar\varepsilon$ & erreur d'approximation et sa borne & N$\cdot$m & --\\
$\Gamma$, $\sigma$ & gain d'adaptation, $\sigma$-modification & --, s$^{-1}$ & $10^7$, 0.02\\
$V$ & fonction de Lyapunov & J & --\\
$\beta$, $c$, $\delta(t)$ & constantes et terme de la preuve (Théorème 1) & -- & $c=4$~s$^{-1}$\\
\bottomrule
\end{longtable}

%======================================================================
\section{Conventions utilisées dans l'article}
\begin{itemize}[itemsep=3pt]
\item \textbf{Convention récepteur (moteur).} Les tensions et courants de la machine sont orientés comme pour un
moteur. En fonctionnement générateur, le couple $T_e$ et le courant $i_{sq}$ sont donc \textbf{négatifs}, et la
puissance électrique fournie par la machine est $P_{dc}=-\frac32\bv v_s^T\bv i_s>0$. C'est la convention
de l'article de référence d'Aimene et al. [8], et c'est sur ce point que les signes avaient été corrigés.
\item \textbf{Transformation de Park à amplitude constante.} Les grandeurs $dq$ ont l'amplitude des
grandeurs de phase (crête). C'est pourquoi les puissances ont un facteur $\frac32$ : $P=\frac32(v_di_d+v_qi_q)$.
\item \textbf{Repère.} Pour la machine, l'axe $d$ est aligné sur le flux des aimants ; pour le réseau,
l'axe $d$ est aligné sur la tension du réseau ($v_{gq}=0$), ce qui suppose une PLL idéale.
\item \textbf{Erreurs.} Toutes les erreurs sont définies comme \og mesure moins référence\fg{} :
$e_\Omega=\Omega-\Omega^*$, $\bv e_s=\bv i_s-\bv i_s^*$, etc.
\item \textbf{Notation vectorielle.} $\bv i_s=[i_{sd}\;i_{sq}]^T$ ; $\|\bv x\|^2=x_d^2+x_q^2$ ;
$\bv x^T\bv y=x_dy_d+x_qy_q$. La matrice $\Jm=\bigl[\begin{smallmatrix}0&1\\-1&0\end{smallmatrix}\bigr]$
fait tourner un vecteur de $-90^\circ$ : $\Jm\bv i=[i_q\;\;-i_d]^T$. Propriété clé :
$\bv x^T\Jm\bv x=x_dx_q-x_qx_d=0$ pour tout $\bv x$. Physiquement, un terme en $\Jm$ échange de l'énergie
entre les axes sans en créer ni en dissiper : c'est un terme \og sans travail\fg{} (\emph{workless}).
\item \textbf{Points.} $\dot x=dx/dt$ ; $\ddot x=d^2x/dt^2$ ; l'étoile $^*$ désigne une référence ;
le chapeau $\hat{\ }$ une estimation ; le tilde $\tilde{\ }$ une erreur d'estimation.
\end{itemize}

%======================================================================
\section{Explication des équations}
La numérotation est la même que dans l'article.

\subsection{Turbine (Section 2.2 de l'article)}
\begin{equation}\tag{1}
\lambda=\frac{\Omega R}{v},\qquad P_w=\tfrac{1}{2}\rho\pi R^2 C_p(\lambda,\beta)\,v^3,\qquad T_w=\frac{P_w}{\Omega}
\end{equation}
\begin{expl}
\sens $\lambda$ compare la vitesse du bout de pale ($\Omega R$) à la vitesse du vent. $\frac12\rho\pi R^2v^3$
est la puissance cinétique du vent qui traverse le disque du rotor ; la turbine n'en capte que la fraction
$C_p$. $T_w$ est le couple que le vent applique sur l'arbre.
\termes $\pi R^2$ : surface balayée ; $v^3$ : la puissance varie avec le cube du vent (vent $\times2$ $\Rightarrow$
puissance $\times8$).
\role c'est la source de puissance du système, et c'est le terme non mesuré que le réseau RBF aide à estimer.
\end{expl}

\begin{equation}\tag{2}
C_p = 0.5176\Bigl(\frac{116}{\lambda_i}-0.4\beta-5\Bigr)e^{-21/\lambda_i}+0.0068\lambda,\qquad
\frac{1}{\lambda_i}=\frac{1}{\lambda+0.08\beta}-\frac{0.035}{\beta^3+1}
\end{equation}
\begin{expl}
\sens formule empirique de Heier [29] qui donne la courbe $C_p(\lambda)$ en forme de cloche. Pour $\beta=0$, le maximum est $C_{p,\max}=0.48$ à $\lambda^\star=8.1$.
\origine ajustement de mesures sur des turbines réelles ; c'est le modèle le plus utilisé dans la littérature.
\num vent nominal : $P_n=\frac12\rho\pi R^2C_{p,\max}v_n^3$ $\Rightarrow$
$v_n=\bigl(2P_n/(\rho\pi R^2C_{p,\max})\bigr)^{1/3}=\bigl(4\cdot10^6/2956\bigr)^{1/3}=11.0$~m/s ;
vitesse nominale $\Omega_n=\lambda^\star v_n/R=8.1\times11.0/40=2.23$~rad/s (21.3~tr/min) ;
couple nominal $T_n=P_n/\Omega_n=0.897$~MN$\cdot$m.
\end{expl}

\subsection{Générateur PMSG (Section 2.3)}
\begin{equation}\tag{3}
v_{sd} = R_s i_{sd} + L_s\frac{di_{sd}}{dt} - \omega_e L_s i_{sq}
\end{equation}
\begin{equation}\tag{4}
v_{sq} = R_s i_{sq} + L_s\frac{di_{sq}}{dt} + \omega_e L_s i_{sd} + \omega_e\psi_f
\end{equation}
\begin{expl}
\sens loi des mailles du stator dans le repère tournant : tension $=$ chute résistive $+$ tension
inductive $+$ couplage $+$ force électromotrice.
\termes $R_si$ : pertes Joule ; $L_s\,di/dt$ : variation du courant ; $\mp\omega_eL_si_{q,d}$ :
termes de couplage entre les axes $d$ et $q$ (ils viennent de la rotation du repère) ;
$\omega_e\psi_f$ : force électromotrice (FEM) des aimants, seulement sur l'axe $q$.
\num à $\Omega_n$ : $\omega_e=26\times2.23=58$~rad/s ($f_e=9.2$~Hz), FEM $\omega_e\psi_f=532$~V,
proche de la tension crête du réseau (563~V) ; réactance $\omega_eL_s=0.091~\Omega$.
\end{expl}

\begin{equation}\tag{5}
T_e = \tfrac32 p\,\psi_f\, i_{sq}
\end{equation}
\begin{expl}
\sens pour une machine à pôles lisses ($L_d=L_q$), le couple ne dépend que de $i_{sq}$. C'est pourquoi on
impose $i_{sd}^*=0$ : $i_{sd}$ ne produit pas de couple et ne ferait qu'ajouter des pertes.
\num $\kappa=\frac32p\psi_f=358$~N$\cdot$m/A ; au couple nominal, $|i_{sq}|=897\,000/358=2.5$~kA.
En générateur, $T_e<0$ donc $i_{sq}<0$.
\end{expl}

\begin{equation}\tag{6}
L_s\dot{\bv{i}}_s = \bv{v}_s - R_s\bv{i}_s + \omega_eL_s\Jm\bv{i}_s - \omega_e\psi_f\bv{e}_q
\end{equation}
\begin{expl}
\sens c'est (3)--(4) écrites sous forme vectorielle, en isolant la dérivée du courant. En développant :
$\omega_eL_s\Jm\bv i_s=\omega_eL_s[i_{sq}\;\;-i_{sd}]^T$, on retrouve exactement $+\omega_eL_si_{sq}$ dans
l'équation $d$ et $-\omega_eL_si_{sd}$ dans l'équation $q$.
\role cette forme fait apparaître la matrice antisymétrique $\Jm$, qui est la clé de la commande par
passivité (Section 5) et du bilan d'énergie (10).
\end{expl}

\subsection{Arbre mécanique (Section 2.4)}
\begin{equation}\tag{7}
J\frac{d\Omega}{dt} = T_w + T_e - B\Omega
\end{equation}
\begin{expl}
\sens principe fondamental de la dynamique en rotation : le couple du vent accélère le rotor, le couple
du générateur ($T_e<0$) le freine, le frottement $B\Omega$ aussi. Si $T_w=-T_e+B\Omega$, la vitesse est
constante.
\num énergie cinétique à $\Omega_n$ : $\frac12J\Omega_n^2=6.96$~MJ ; constante d'inertie
$H=6.96/2=3.5$~s (valeur typique des grandes éoliennes). Cette grande inertie explique pourquoi la
vitesse ne peut pas suivre instantanément les variations rapides du vent.
\end{expl}

\subsection{Bus continu et côté réseau (Section 2.5)}
\begin{equation}\tag{8}
CV_{dc}\frac{dV_{dc}}{dt} = P_{dc} - P_i,\qquad P_{dc}=-\tfrac32\bv{v}_s^T\bv{i}_s,\qquad
P_i=\tfrac32\bv{v}_i^T\bv{i}_g
\end{equation}
\begin{expl}
\sens l'énergie du condensateur $\frac12CV_{dc}^2$ a pour dérivée $CV_{dc}\dot V_{dc}$. Elle augmente si la
puissance qui arrive du MSC ($P_{dc}$) est supérieure à celle qui part vers le GSC ($P_i$). Le signe $-$ de
$P_{dc}$ vient de la convention récepteur : la puissance absorbée par la machine est
$\frac32\bv v_s^T\bv i_s$, donc la puissance fournie est son opposé.
\num énergie du bus : $\frac12\times0.02\times1200^2=14.4$~kJ, soit seulement 7~ms de puissance
nominale. D'où la nécessité d'une boucle de bus rapide et d'une anticipation de $P_{dc}$.
\end{expl}

\begin{equation}\tag{9}
L_f\dot{\bv{i}}_g = \bv{v}_i - R_f\bv{i}_g + \omega_gL_f\Jm\bv{i}_g - \bv{v}_g
\end{equation}
\begin{expl}
\sens loi des mailles du filtre entre la tension du GSC $\bv v_i$ et la tension du réseau $\bv v_g$, dans
le repère du réseau. Même structure que (6) : résistance, inductance, couplage en $\Jm$.
\termes $P_g=\frac32v_{gd}i_{gd}$ et $Q_g=-\frac32v_{gd}i_{gq}$ : avec $v_{gq}=0$, la puissance active ne
dépend que de $i_{gd}$ et la réactive que de $i_{gq}$ ; $i_{gq}^*=0$ donne $Q_g=0$.
\num $v_{gd}=690\sqrt{2/3}=563$~V ; courant nominal crête $I_n=2P_n/(3\times563)=2.37$~kA.
\end{expl}

\subsection{Bilan d'énergie (Section 2.6)}
\begin{equation}\tag{10}
\dot E = P_w - B\Omega^2 - \tfrac32R_s\|\bv{i}_s\|^2 - \tfrac32R_f\|\bv{i}_g\|^2 - P_g
\end{equation}
\begin{expl}
\sens la puissance du vent est soit stockée ($\dot E$), soit perdue (frottement, résistances), soit injectée
au réseau. C'est la conservation de l'énergie.
\origine on dérive $E=\frac12J\Omega^2+\frac34L_s\|\bv i_s\|^2+\frac12CV_{dc}^2+\frac34L_f\|\bv i_g\|^2$ et on
remplace avec (6)--(9) :
\begin{itemize}[itemsep=1pt]
\item arbre : $J\Omega\dot\Omega=T_w\Omega+T_e\Omega-B\Omega^2$ ;
\item machine : $\frac32\bv i_s^TL_s\dot{\bv i}_s=\frac32\bv i_s^T\bv v_s-\frac32R_s\|\bv i_s\|^2-\frac32\omega_e\psi_fi_{sq}$
(le terme en $\Jm$ disparaît car $\bv i_s^T\Jm\bv i_s=0$) ;
\item bus : $CV_{dc}\dot V_{dc}=-\frac32\bv v_s^T\bv i_s-\frac32\bv v_i^T\bv i_g$ ;
\item filtre : $\frac32\bv i_g^TL_f\dot{\bv i}_g=\frac32\bv i_g^T\bv v_i-\frac32R_f\|\bv i_g\|^2-\frac32\bv i_g^T\bv v_g$.
\end{itemize}
En additionnant, les puissances des convertisseurs s'annulent deux à deux, et
$T_e\Omega=\frac32p\psi_fi_{sq}\Omega=\frac32\omega_e\psi_fi_{sq}$ annule le terme de FEM. Il reste (10),
avec $P_g=\frac32\bv i_g^T\bv v_g$.
\role c'est la \textbf{preuve que tous les signes du modèle sont cohérents} (réponse directe à la remarque du
professeur sur les signes). Numériquement, le résidu est inférieur à $10^{-15}P_n$.
\end{expl}

\subsection{Objectifs et MPPT (Section 3)}
\begin{equation}\tag{11}
\Omega_{opt}=\frac{\lambda^\star v_m}{R}
\end{equation}
\begin{expl}
\sens pour capter le maximum de puissance, il faut $\lambda=\lambda^\star$, donc $\Omega=\lambda^\star v/R$
(MPPT de type TSR). On utilise le vent mesuré $v_m$, limité à $[0.5\Omega_n,\Omega_n]$.
\num à 9~m/s : $\Omega_{opt}=8.1\times9/40=1.82$~rad/s ; à 11~m/s : 2.23~rad/s.
\end{expl}

\begin{equation}\tag{12}
\ddot\Omega^*=\omega_r^2(\Omega_{opt}-\Omega^*)-2\omega_r\dot\Omega^*
\end{equation}
\begin{expl}
\sens filtre du second ordre de fonction de transfert $\Omega^*/\Omega_{opt}=\omega_r^2/(s+\omega_r)^2$
(amortissement critique, sans dépassement). Il donne une référence lisse $\Omega^*$ \textbf{et sa dérivée}
$\dot\Omega^*$, nécessaire à la loi plate (14).
\role sans ce filtre, un échelon de vent demanderait un échelon de vitesse, impossible avec
$J=2.8\times10^6$~kg$\cdot$m$^2$ (couple infini). Avec $\omega_r=1$~rad/s, la référence atteint sa nouvelle
valeur en 4 à 5~s.
\end{expl}

\subsection{Platitude (Section 4)}
\textbf{Définition.} Un système est \emph{plat} s'il existe une sortie $y$ (sortie plate) telle que tous les
états et toutes les entrées s'écrivent en fonction de $y$ et d'un nombre fini de ses dérivées [17]. On
peut alors calculer directement l'entrée qui produit une trajectoire désirée $y^*(t)$.

\begin{equation}\tag{13}
T_e=J\dot\Omega+B\Omega-T_w,\qquad i_{sq}=\frac{T_e}{\frac32p\psi_f}
\end{equation}
\begin{expl}
\sens c'est (7) et (5) inversées : si l'on connaît $\Omega(t)$ et $\dot\Omega(t)$, on connaît le couple et le
courant nécessaires. Avec $i_{sd}$, puis (3)--(4), on obtient aussi les tensions : $(\Omega,i_{sd})$ est une
sortie plate du système PMSG $+$ arbre (Table 2).
\end{expl}

\begin{equation}\tag{14}
\nu_1=\dot\Omega^*-k_{11}e_\Omega-k_{12}\!\int\!e_\Omega\,dt,\qquad e_\Omega=\Omega-\Omega^*
\end{equation}
\begin{equation}\tag{15}
T_e^* = J\nu_1 + B\Omega - \hat T_w - \hat\Delta,\qquad i_{sq}^*=\frac{T_e^*}{\frac32p\psi_f},\qquad i_{sd}^*=0
\end{equation}
\begin{expl}
\sens on remplace $\dot\Omega$ dans (13) par une accélération \og désirée\fg{} $\nu_1$ : la dérivée de la
référence plus une correction proportionnelle et intégrale de l'erreur. $\hat T_w$ est le couple du vent
calculé par le modèle (1)--(2) avec $v_m$ ; $\hat\Delta$ est la correction du réseau RBF.
\origine en remplaçant (15) dans (7), avec $T_e=T_e^*$ (courant bien suivi) :
$J\dot\Omega=T_w+(J\nu_1+B\Omega-\hat T_w-\hat\Delta)-B\Omega=J\nu_1+(T_w-\hat T_w)-\hat\Delta$.
Comme $\nu_1=\dot\Omega^*-k_{11}e_\Omega-k_{12}\int e_\Omega$, on obtient l'équation d'erreur (23) ci-dessous.
\num gains : l'équation d'erreur sans incertitude est $\ddot e+k_{11}\dot e+k_{12}e=0$ (en dérivant), de
polynôme $s^2+k_{11}s+k_{12}$. Pour un pôle double en $-\omega_\Omega$ : $(s+\omega_\Omega)^2=s^2+2\omega_\Omega s+\omega_\Omega^2$,
donc $k_{11}=2\omega_\Omega=4$, $k_{12}=\omega_\Omega^2=4$ avec $\omega_\Omega=2$~rad/s.
\role $T_e^*$ est limité à $[-1.2T_n,0]$ (générateur uniquement) ; l'intégrale est gelée à la limite
(anti-emballement).
\end{expl}

\begin{equation}\tag{16}
\dot y_2 = P_{dc} - P_g - \tfrac32R_f\|\bv{i}_g\|^2,\qquad y_2=\tfrac12CV_{dc}^2+\tfrac34L_f\|\bv i_g\|^2
\end{equation}
\begin{expl}
\sens $y_2$ est l'énergie stockée dans le bus \textbf{et} dans le filtre. Sa dérivée est une simple balance
de puissances : ce qui arrive ($P_{dc}$) moins ce qui part ($P_g$) moins les pertes du filtre.
\origine en dérivant $y_2$ avec (8) et (9) : les termes $\pm\frac32\bv v_i^T\bv i_g$ s'annulent et le
terme en $\Jm$ disparaît. C'est pour cela qu'on ajoute l'énergie du filtre : l'équation devient
\textbf{linéaire en $P_g$}, sans dérivée de courant [18].
\end{expl}

\begin{equation}\tag{17}
P_g^* = P_{dc} - \tfrac32R_f\|\bv{i}_g\|^2 + k_{21}e_y + k_{22}\!\int\! e_y\,dt,\qquad
i_{gd}^*=\frac{2P_g^*}{3v_{gd}},\qquad i_{gq}^*=0
\end{equation}
\begin{expl}
\sens on impose $\dot y_2=-k_{21}e_y-k_{22}\int e_y$ ; en remplaçant dans (16) on trouve la puissance à
injecter. Le terme $P_{dc}$ est une \textbf{anticipation} : la puissance du générateur est envoyée au réseau
avant qu'elle ne fasse monter $V_{dc}$.
\num pôle double en $-\omega_v$ : $k_{21}=2\omega_v=120$, $k_{22}=\omega_v^2=3600$ avec
$\omega_v=60$~rad/s. $i_{gd}^*=2P_g^*/(3v_{gd})$ vient de $P_g=\frac32v_{gd}i_{gd}$.
\role dans la simulation réaliste, $P_{dc}$ et $v_{gd}$ sont filtrés (2~ms) pour ne pas transmettre le bruit
de mesure à la référence.
\end{expl}

\subsection{Passivité (Section 5)}
\textbf{Idée.} Au lieu d'annuler toutes les non-linéarités, la PBC (i) impose une énergie désirée à la
boucle fermée (\emph{energy shaping}) et (ii) ajoute de l'amortissement (\emph{damping injection}), en
gardant les termes physiques qui ne produisent pas d'énergie (termes en $\Jm$) [19,20].

\begin{equation}\tag{18}
L_s\dot{\bv{i}}_s = \bigl(\omega_eL_s\Jm - R_s\bv{I}\bigr)\bv{i}_s + \bv{v}_s - \omega_e\psi_f\bv{e}_q ,
\qquad \dot H_s = -R_s\|\bv{i}_s\|^2 + \bv{i}_s^T(\bv{v}_s-\omega_e\psi_f\bv{e}_q)
\end{equation}
\begin{expl}
\sens c'est (6) écrite sous forme hamiltonienne : $\omega_eL_s\Jm$ (antisymétrique) $=$ interconnexion,
$R_s\bv I$ (symétrique positive) $=$ dissipation. Avec l'énergie magnétique $H_s=\frac12L_s\|\bv i_s\|^2$,
la dérivée montre que l'énergie varie seulement par la dissipation $-R_s\|\bv i_s\|^2$ et par la puissance
fournie au port $\bv i_s^T(\ldots)$ : le système est \textbf{passif}.
\end{expl}

\begin{equation}\tag{19}
\bv{v}_s^* = \underbrace{L_s\dot{\bv{i}}_s^* + R_s\bv{i}_s^* - \omega_eL_s\Jm\bv{i}_s^* + \omega_e\psi_f\bv{e}_q}_{\text{mise en forme de l'énergie}}
\underbrace{- R_a\bv{e}_s}_{\text{amortissement}} - K_i\bv{z}_s
\end{equation}
\begin{expl}
\sens la première partie est le modèle (6) appliqué à la référence : c'est la tension qui produirait
exactement $\bv i_s^*$. $-R_a\bv e_s$ ajoute une résistance virtuelle $R_a$ qui amortit l'erreur ;
$-K_i\bv z_s$ ($\dot{\bv z}_s=\bv e_s$) supprime l'erreur statique due aux erreurs de paramètres.
\termes $\dot{\bv i}_s^*$ est obtenue par un filtre du premier ordre de 2~ms.
\end{expl}

\begin{equation}\tag{20}
L_s\dot{\bv{e}}_s=-(R_s+R_a)\bv{e}_s+\omega_eL_s\Jm\bv{e}_s-K_i\bv{z}_s
\end{equation}
\begin{expl}
\origine on soustrait $L_s\dot{\bv i}_s^*$ de (6) et on remplace $\bv v_s$ par (19) : la FEM
$\omega_e\psi_f\bv e_q$ s'annule, et il reste une dynamique d'erreur qui a la \textbf{même structure} que la
machine (couplage en $\Jm$ conservé) avec une résistance totale $R_s+R_a$.
\num pôle double en $-\omega_i$ par axe : $L_ss^2+(R_s+R_a)s+K_i=L_s(s+\omega_i)^2$, d'où
$R_a=2\omega_iL_s-R_s=2\times500\times1.57\cdot10^{-3}-0.78\cdot10^{-3}=1.57~\Omega$ et
$K_i=\omega_i^2L_s=392.5~\Omega$/s.
\end{expl}

\begin{equation}\tag{21}
\bv{v}_i^* = L_f\dot{\bv{i}}_g^* + R_f\bv{i}_g^* - \omega_gL_f\Jm\bv{i}_g^* + \bv{v}_g
- R_b\bv{e}_g - K_{ig}\bv{z}_g
\end{equation}
\begin{expl}
\sens même loi pour le filtre du réseau ; $+\bv v_g$ compense la tension du réseau.
\num $R_b=2\times500\times0.15\cdot10^{-3}-2\cdot10^{-3}=0.148~\Omega$, $K_{ig}=500^2\times0.15\cdot10^{-3}=37.5~\Omega$/s.
\end{expl}

\textbf{Proposition 1 (expliquée).} Avec $H_e=\frac12L_s\|\bv e_s\|^2+\frac12K_i\|\bv z_s\|^2$ (énergie
\og magnétique\fg{} de l'erreur $+$ énergie stockée dans l'intégrateur) :
$\dot H_e=\bv e_s^TL_s\dot{\bv e}_s+K_i\bv z_s^T\dot{\bv z}_s$. En remplaçant (20) et $\dot{\bv z}_s=\bv e_s$ :
$\dot H_e=-(R_s+R_a)\|\bv e_s\|^2+\omega_eL_s\underbrace{\bv e_s^T\Jm\bv e_s}_{=0}-K_i\bv e_s^T\bv z_s+K_i\bv z_s^T\bv e_s
=-(R_s+R_a)\|\bv e_s\|^2\le0$. L'énergie de l'erreur ne peut que diminuer : la boucle de courant est
stable. Les gains choisis séparent les bandes passantes d'une décade : courant 500, bus 60, vitesse
2~rad/s.

\subsection{Réseau RBF (Section 6)}
\begin{equation}\tag{22}
\Delta(\bv{x},t)=T_w-\hat T_w
\end{equation}
\begin{expl}
\sens l'erreur entre le vrai couple du vent et le couple calculé par le modèle. Elle regroupe l'erreur du
modèle $C_p$, le retard de l'anémomètre (qui mesure le vent en un point, avec 1~s de retard) et l'erreur
sur les paramètres du générateur (par exemple $\psi_f$).
\num dans le cas 3, sa moyenne vaut 83~kN$\cdot$m, soit 14\,\% du couple du vent.
\end{expl}

\begin{equation}\tag{23}
J\bigl(\dot e_\Omega + k_{11}e_\Omega + k_{12}\textstyle\int e_\Omega\,dt\bigr) = \Delta-\hat\Delta
\end{equation}
\begin{expl}
\sens équation de l'erreur de vitesse obtenue en mettant (15) dans (7). Si $\hat\Delta=\Delta$, le second
membre est nul et l'erreur tend vers zéro avec les pôles choisis. Sinon, $\Delta-\hat\Delta$ agit comme une
perturbation : le but du réseau est de rendre ce terme petit.
\end{expl}

\begin{equation}\tag{24}
\bv{z}=\Bigl[\frac{\Omega}{\Omega_n}\;\;\frac{v_m}{v_n}\Bigr]^T
\end{equation}
\begin{expl}
\sens l'entrée du réseau : le point de fonctionnement (vitesse et vent), normalisé pour être de l'ordre
de 1. $\Delta$ dépend surtout de ce point de fonctionnement.
\end{expl}

\begin{equation}\tag{25}
\Delta=\bv{W}^{*T}\boldsymbol\Phi(\bv{z})+\varepsilon,\qquad |\varepsilon|\le\bar\varepsilon ,\qquad
\Phi_j(\bv{z})=\exp\Bigl(-\frac{\|\bv{z}-\bv{c}_j\|^2}{2b^2}\Bigr)
\end{equation}
\begin{expl}
\sens théorème d'approximation universelle [21] : sur une région bornée, une fonction continue peut être
approchée par une somme pondérée de gaussiennes, avec une erreur $\varepsilon$ aussi petite que l'on veut si
l'on met assez de neurones. $\bv W^*$ sont les \og meilleurs\fg{} poids (inconnus).
\termes chaque gaussienne $\Phi_j$ vaut 1 au centre $\bv c_j$ et décroît avec la distance ; les 25 centres
forment une grille $5\times5$ avec les valeurs 0.5, 0.65, 0.8, 0.95 et 1.1 sur chaque axe ; la largeur
$b=0.15$ est égale à l'espacement de la grille, donc les gaussiennes se recouvrent.
\end{expl}

\begin{equation}\tag{26}
\hat\Delta(\bv{z})=\hat{\bv{W}}^T\boldsymbol\Phi(\bv{z})
\end{equation}
\begin{expl}
\sens la sortie du réseau, avec les poids estimés $\hat{\bv W}$.
\role avec le réseau, $k_{12}=0$ : sinon l'intégrale de la boucle de vitesse absorberait la partie lente de
$\Delta$, l'erreur $e_\Omega$ resterait nulle en moyenne et le réseau n'apprendrait rien (constaté en
simulation).
\end{expl}

\begin{equation}\tag{27}
\dot{\hat{\bv{W}}} = \Gamma\bigl(\boldsymbol\Phi(\bv{z})\,e_\Omega - \sigma\hat{\bv{W}}\bigr),\qquad \hat{\bv{W}}(0)=0
\end{equation}
\begin{expl}
\sens loi d'apprentissage en ligne : chaque poids $j$ change proportionnellement à l'erreur de vitesse et à
l'activation de son neurone. Seuls les neurones proches du point de fonctionnement actuel apprennent.
\termes $\Gamma$ : vitesse d'apprentissage ; $-\sigma\hat{\bv W}$ ($\sigma$-modification [23]) : un léger
\og oubli\fg{} qui empêche les poids de dériver à cause du bruit ; sans lui, les poids pourraient croître
indéfiniment. L'adaptation est gelée quand $T_e^*$ est saturé.
\num $\Gamma=10^7$ paraît grand, mais il faut le comparer aux unités : $e_\Omega\approx10^{-2}$~rad/s et
$\Phi\le1$, donc $\dot{\hat W}\approx10^5$~N$\cdot$m/s : le réseau apprend une correction de l'ordre de
100~kN$\cdot$m en environ 1~s.
\origine cette loi est choisie pour annuler le terme croisé dans la dérivée de la fonction de Lyapunov
(voir la preuve, équation (32)).
\end{expl}

\subsection{Stabilité (Section 7)}
\textbf{Les hypothèses.} (H1) L'approximation (25) est valable sur la région de fonctionnement. (H2) $J$ et
$B$ sont connus, les erreurs de paramètres sont incluses dans $\Delta$, il n'y a pas de saturation, et les
références et la tension du réseau sont bornées. Ce sont des hypothèses classiques ; la Remarque 1 de
l'article dit honnêtement que l'échantillonnage, les saturations et le bruit ne sont pas couverts par la
preuve, mais vérifiés par simulation.

\textbf{Le théorème en mots.} (i) Les erreurs de courant tendent vers zéro. (ii) L'erreur de vitesse et
l'erreur des poids finissent dans une petite région autour de zéro, dont la taille dépend de l'erreur
d'approximation $\bar\varepsilon$ et de $\sigma$. (iii) L'erreur d'énergie du bus reste bornée. Donc tous
les signaux restent bornés.

\begin{equation}\tag{28}
\limsup_{t\to\infty}\Bigl(\tfrac12Je_\Omega^2+\tfrac{1}{2\Gamma}\|\tilde{\bv{W}}\|^2\Bigr)\le\frac{\beta}{c},
\qquad \beta=\frac{\bar\varepsilon^2}{Jk_{11}}+\frac{\sigma\bar W^2}{2},\qquad c=\min(k_{11},\sigma\Gamma)
\end{equation}
\begin{expl}
\sens c'est la taille de la région finale (borne UUB). $\limsup$ = la plus grande valeur que la quantité
prend \og à la fin\fg{}. Plus $\bar\varepsilon$ est petit, plus la borne est petite.
\num $c=\min(4,\;0.02\times10^7)=4$~s$^{-1}$ : la convergence vers la région finale a une constante de
temps d'environ 0.25~s pour cette fonction d'énergie.
\end{expl}

\begin{equation}\tag{29}
\int_0^\infty\|\bv{e}_s\|^2dt\le\frac{H_e(0)}{R_s+R_a}
\end{equation}
\begin{expl}
\origine on intègre $\dot H_e=-(R_s+R_a)\|\bv e_s\|^2$ de 0 à $t$ : $H_e(t)-H_e(0)=-(R_s+R_a)\int_0^t\|\bv e_s\|^2$,
et comme $H_e(t)\ge0$, on obtient (29). L'erreur de courant a une \og énergie totale\fg{} finie.
\role avec le lemme de Barbalat [30] (si une fonction a une intégrale finie et une dérivée bornée, elle
tend vers zéro), on conclut $\bv e_s\to0$.
\end{expl}

\begin{equation}\tag{30}
J\dot e_\Omega=-Jk_{11}e_\Omega+\tilde{\bv{W}}^T\boldsymbol\Phi+\varepsilon+\kappa e_{sq}
\end{equation}
\begin{expl}
\origine (23) avec $k_{12}=0$, en remplaçant $\Delta-\hat\Delta=\bv W^{*T}\boldsymbol\Phi+\varepsilon-\hat{\bv W}^T\boldsymbol\Phi
=\tilde{\bv W}^T\boldsymbol\Phi+\varepsilon$, et en tenant compte que le courant n'est pas exactement égal à
sa référence : $T_e=\kappa i_{sq}=\kappa(i_{sq}^*+e_{sq})=T_e^*+\kappa e_{sq}$. C'est le terme qui
\textbf{couple} la boucle de courant et la boucle de vitesse.
\end{expl}

\begin{equation}\tag{31}
V=\tfrac12Je_\Omega^2+\tfrac{1}{2\Gamma}\|\tilde{\bv{W}}\|^2
\end{equation}
\begin{expl}
\sens fonction de Lyapunov : \og énergie cinétique\fg{} de l'erreur de vitesse $+$ \og énergie\fg{} de l'erreur
des poids. Elle est positive et nulle seulement si les deux erreurs sont nulles.
\end{expl}

\begin{equation}\tag{32}
\dot V=-Jk_{11}e_\Omega^2+e_\Omega\varepsilon+\kappa e_\Omega e_{sq}+\sigma\tilde{\bv{W}}^T\hat{\bv{W}}
\end{equation}
\begin{expl}
\origine $\dot V=e_\Omega J\dot e_\Omega+\frac1\Gamma\tilde{\bv W}^T\dot{\tilde{\bv W}}$, avec
$\dot{\tilde{\bv W}}=-\dot{\hat{\bv W}}$ (car $\bv W^*$ est constant). En remplaçant (30) et (27) :
le terme $e_\Omega\tilde{\bv W}^T\boldsymbol\Phi$ venant de (30) et le terme $-\tilde{\bv W}^T\boldsymbol\Phi e_\Omega$
venant de (27) \textbf{s'annulent} : c'est exactement pour cela que la loi (27) a cette forme. Il reste (32).
\end{expl}

\begin{equation}\tag{33}
\dot V\le-\frac{Jk_{11}}{2}e_\Omega^2-\frac\sigma2\|\tilde{\bv{W}}\|^2+\beta+\delta(t)\le-cV+\beta+\delta(t),
\qquad \delta(t)=\frac{\kappa^2}{Jk_{11}}e_{sq}^2(t)
\end{equation}
\begin{expl}
\origine trois majorations classiques :
\begin{itemize}[itemsep=1pt]
\item inégalité de Young $ab\le\frac{\gamma}{2}a^2+\frac{1}{2\gamma}b^2$ :
$e_\Omega\varepsilon\le\frac{Jk_{11}}{4}e_\Omega^2+\frac{\bar\varepsilon^2}{Jk_{11}}$ et
$\kappa e_\Omega e_{sq}\le\frac{Jk_{11}}{4}e_\Omega^2+\frac{\kappa^2}{Jk_{11}}e_{sq}^2$ ;
\item identité $\tilde{\bv W}^T\hat{\bv W}=\tilde{\bv W}^T(\bv W^*-\tilde{\bv W})\le\frac12\|\bv W^*\|^2-\frac12\|\tilde{\bv W}\|^2$ ;
\item comme $\frac{Jk_{11}}{2}e_\Omega^2=k_{11}\cdot\frac12Je_\Omega^2$ et
$\frac\sigma2\|\tilde{\bv W}\|^2=\sigma\Gamma\cdot\frac{1}{2\Gamma}\|\tilde{\bv W}\|^2$, ces deux termes sont $\ge cV$.
\end{itemize}
\role le lemme de comparaison donne $V(t)\le V(0)e^{-ct}+\beta/c+\int_0^te^{-c(t-\tau)}\delta(\tau)d\tau$ :
le premier terme disparaît, le dernier aussi (car $\delta$ a une intégrale finie grâce à (29)), et il reste
$\beta/c$, c'est-à-dire (28).
\end{expl}

\textbf{Point (iii) du théorème.} Avec $\eta=\int e_y\,dt$, les équations (16)--(17) donnent
$\ddot\eta+k_{21}\dot\eta+k_{22}\eta=-\frac32v_{gd}e_{gd}-\dot y_2^*$. Le polynôme $s^2+120s+3600=(s+60)^2$
est stable (Hurwitz : toutes ses racines sont à partie réelle négative), l'entrée $e_{gd}$ tend vers zéro
et $\dot y_2^*$ (variation de l'énergie de référence du filtre) est petite et bornée : l'erreur reste bornée.

\subsection{Simulation (Section 8)}
\begin{equation}\tag{34}
v(t) = V_{mean} + \sum_{k=1}^{3} A_k\sin(2\pi f_k t) + v_{turb}(t)
\end{equation}
\begin{expl}
\sens vent aléatoire : une moyenne de 8.5~m/s, trois sinusoïdes lentes (périodes 33, 6.3 et 5.9~s), comme
dans l'article de Mathematics [27], plus une turbulence $v_{turb}$.
\termes $v_{turb}$ est un bruit coloré du premier ordre :
$v_{turb}[k+1]=a\,v_{turb}[k]+\sqrt{1-a^2}\,\sigma_t\,w[k]$, $a=e^{-\Delta t/\tau_t}$, avec
$\sigma_t=0.5$~m/s, $\tau_t=1.5$~s, $\Delta t=1$~ms et $w$ un bruit blanc gaussien (graine fixe, donc
résultats reproductibles).
\end{expl}

\begin{equation}\tag{35}
\begin{gathered}
\mathrm{RMSE}=\sqrt{\frac1N\sum_{k=1}^Ne_\Omega^2(k)},\qquad
\mathrm{MAE}=\frac1N\sum_{k=1}^N|e_\Omega(k)|,\qquad
e_{\max}=\max_k|e_\Omega(k)|,\\
\mathrm{IAE}=\int|e_\Omega|\,dt,\qquad
\mathrm{ISE}=\int e_\Omega^2\,dt
\end{gathered}
\end{equation}
\begin{expl}
\sens cinq mesures de l'erreur de vitesse calculées de 1 à 40~s : RMSE et ISE pénalisent les grandes
erreurs, MAE et IAE l'erreur moyenne, $e_{\max}$ le pire instant. IAE et MAE donnent le même pourcentage
d'amélioration car IAE $=$ MAE $\times$ durée.
\end{expl}

\textbf{Autres indicateurs des Tables 4 et 5.}
\begin{itemize}[itemsep=2pt]
\item \emph{Temps d'établissement} : temps après l'échelon de vent pour que $|e_\Omega|$ reste sous 2~mrad/s.
\item \emph{Rendement MPPT} : $\eta=\int P_w\,dt\big/\int P_{\max}\,dt$, avec $P_{\max}=\frac12\rho\pi R^2C_{p,\max}v^3$.
\item \emph{Ondulation de couple} : écart type de $T_e$ moins sa moyenne glissante sur 50~ms, en \% de $T_n$.
\item \emph{Activité du couple} : écart type de la dérivée du couple moyenné sur 50~ms (MN$\cdot$m/s).
\item \emph{THD} : $\mathrm{THD}=\sqrt{\sum_{h=2}^{20}I_h^2}\,/\,I_1$, calculé par FFT sur 10 périodes (0.2~s).
La norme IEEE 519 [31] limite (pour un rapport $I_{cc}/I_L<20$) chaque harmonique impair d'ordre inférieur à 11 à 4\,\%, les harmoniques pairs au quart de cette valeur, et la distorsion totale à 5\,\%.
\end{itemize}

\textbf{Conditions réalistes de la simulation.} Bruit blanc sur les mesures : courants (écart type 5~A,
soit 0.2\,\% du courant nominal), tension du bus et du réseau (1~V), vitesse du codeur (2~mrad/s, filtrée
5~ms). Harmoniques du réseau : rang 5 (2\,\%, séquence inverse) et rang 7 (1.5\,\%, séquence directe). Dans le
repère $dq$, ces deux harmoniques apparaissent toutes les deux à $6\times50=300$~Hz : c'est l'ondulation
visible sur $i_{gq}$ et $P_g$.

%======================================================================
\section{Comment lire les figures et les tables}
\begin{longtable}{@{}p{2.2cm}p{12.6cm}@{}}
\toprule
Élément & Ce qu'il faut regarder\\
\midrule
\endhead
Fig. 1 & Le schéma du système : turbine, PMSG, MSC, bus, GSC, filtre, réseau.\\
Fig. 2 & Comment la platitude calcule $i_{sq}^*$ (en haut) et $i_{gd}^*$ (en bas) à partir des références.\\
Fig. 3 & La PBC : modèle de référence $+$ amortissement $+$ intégrale ; l'énergie $H$ diminue.\\
Fig. 4 & Le contrôleur complet avec les trois couches et le réseau RBF.\\
Fig. 5 & Les deux profils de vent (échelon et aléatoire).\\
Fig. 6 & (a)--(b) la vitesse mesurée suit la référence ; (c) l'erreur est plus petite avec le RBF (rouge) ;
(d) l'IAE cumulée : l'écart entre bleu et rouge grandit avec le temps.\\
Fig. 7 & $C_p$ reste proche de 0.48 ; la puissance captée suit la puissance disponible.\\
Fig. 8 & Les courants mesurés suivent leurs références ; $i_{sd}$ et $i_{gq}$ restent autour de zéro.\\
Fig. 9 & (a) le couple suit sa référence, l'ondulation est faible ; (b) le RBF suit l'incertitude réelle.\\
Fig. 10 & $V_{dc}$ reste à $\pm7$~V ; pendant le creux de tension, le courant double et $V_{dc}$ revient en 50~ms.\\
Fig. 11 & $P_g$ suit le vent, $Q_g\approx0$ (facteur de puissance unitaire), même pendant le creux.\\
Fig. 12 & Le courant du réseau est presque sinusoïdal ; chaque harmonique est sous la limite IEEE 519.\\
Table 1 & Les paramètres et leurs sources (Hasan \& Shatil [28], Rolán [7], Aimene [8]).\\
Table 2 & Les sorties plates et ce qu'on calcule à partir d'elles.\\
Table 3 & Les paramètres du RBF et tous les gains.\\
Table 4 & L'erreur de vitesse sans et avec RBF : amélioration de 26.5 à 44.8\,\% (IAE).\\
Table 5 & Le reste du système ne change pas avec le RBF (bus, courants, THD, énergie) ; seul le couple est
un peu plus actif (9 à 12\,\%).\\
\bottomrule
\end{longtable}

%======================================================================
\section{Questions probables du professeur et réponses}
\begin{enumerate}[itemsep=5pt,leftmargin=*]
\item \textbf{Pourquoi la convention récepteur, et comment être sûr des signes ?}
C'est la convention de l'article de référence [8]. Le bilan d'énergie (10), obtenu analytiquement et
vérifié numériquement (résidu $<10^{-15}P_n$), prouve que chaque signe du modèle est cohérent.

\item \textbf{D'où viennent les valeurs numériques ?}
Le générateur, le bus, le filtre et le réseau viennent de la Table I de Hasan et Shatil [28] ; le rayon de
Rolán et al. [7] ; $\rho$ d'Aimene et al. [8] ; $C_p$ de Heier [29]. $J$, $B$ et $R_f$ sont supposés et
indiqués comme tels ($J$ donne une constante d'inertie de 3.5~s, valeur typique).

\item \textbf{Pourquoi la platitude et pas une simple boucle PI ?}
La platitude calcule directement la commande à partir de la trajectoire désirée : le couple nécessaire
pour suivre $\Omega^*$ et la puissance nécessaire pour garder l'énergie du bus. Pour le bus, l'anticipation
de $P_{dc}$ envoie la puissance au réseau avant que $V_{dc}$ ne bouge.

\item \textbf{Pourquoi l'énergie du filtre est-elle incluse dans $y_2$ ?}
Pour que $\dot y_2$ soit exactement $P_{dc}-P_g-$ pertes, linéaire en $P_g$, sans dérivée du courant.

\item \textbf{Qu'apporte la PBC par rapport à une linéarisation exacte ?}
Elle ne compense pas les termes de couplage en $\Jm$ (ils ne créent pas d'énergie), elle ajoute un
amortissement, et la convergence se démontre simplement avec l'énergie $H_e$ (Proposition 1).

\item \textbf{Pourquoi $k_{12}=0$ avec le réseau de neurones ?}
Sinon l'intégrale absorbe la partie lente de l'incertitude, l'erreur reste nulle en moyenne et le réseau
n'apprend rien. Le réseau remplace l'intégrale.

\item \textbf{Le RBF n'augmente pas l'énergie produite : à quoi sert-il ?}
Il améliore la précision et la robustesse du suivi de vitesse (IAE $-$26.5 à $-$44.8\,\%, temps
d'établissement divisé par deux) sans augmenter les gains. Avec une inertie aussi grande, l'énergie captée
est déjà proche du maximum (97.7\,\%) ; l'article le dit clairement.

\item \textbf{Pourquoi ne pas augmenter simplement $k_{11}$ ?}
Le gain $k_{11}$ est multiplié par $J=2.8\times10^6$ dans (15) : il amplifie le bruit de mesure de la
vitesse. Le réseau réduit l'erreur sans toucher ce gain (Remarque 1).

\item \textbf{La preuve de stabilité couvre-t-elle tout ?}
Non, et l'article le dit : elle suppose $J$, $B$ connus et pas de saturation, et ne couvre pas
l'échantillonnage ni le bruit. Ces effets sont évalués par simulation réaliste.

\item \textbf{Pourquoi un modèle moyen des convertisseurs ?}
Il représente correctement la dynamique de commande (bande passante de 500~rad/s, très inférieure à
$f_{sw}$). L'ondulation de commutation n'est pas représentée ; c'est une limite écrite dans la conclusion,
avec un filtre $LCL$ en perspective.

\item \textbf{La PLL est idéale : est-ce un problème ?}
C'est une hypothèse courante pour étudier la commande ; la dynamique de la PLL est citée comme perspective.

\item \textbf{Pourquoi le THD vaut 3.9\,\% ?}
Parce que la tension du réseau contient 2\,\% d'harmonique 5 et 1.5\,\% d'harmonique 7 ; la boucle de
courant ne les rejette pas totalement. Toutes les harmoniques restent sous les limites IEEE 519.

\item \textbf{Pourquoi n'y a-t-il pas de comparaison avec FOC ou SMC ?}
L'article évalue l'apport du réseau de neurones avec les mêmes gains (sans/avec RBF), comme l'article
publié dans Mathematics [27]. Une comparaison avec d'autres commandes est annoncée comme travail futur.

\item \textbf{Les résultats sont-ils reproductibles ?}
Oui : le script MATLAB est fourni, le vent et le bruit utilisent un générateur avec une graine fixe
(mêmes nombres sous MATLAB et Octave), et le paramètre \texttt{P.real} permet de passer du cas réaliste
au cas idéal.
\end{enumerate}

\end{document}
